\section{Related Work}
\label{sec:related}

\paragraph{Engagement recognition in classroom video.} DAiSEE
~\cite{gupta2016daisee} introduced the standard subject-disjoint
ordinal four-class engagement benchmark (9,068 ten-second clips,
112 students). Early supervised work---LRCN~\cite{donahue2015lrcn},
3D-CNN baselines reported in~\cite{gupta2016daisee}---reaches
$\approx 47\%$ accuracy on the test split. More recent supervised
methods (CavT~\cite{ai2022cavt}, ViBED-Net~\cite{abedi2021vibednet},
EngageNet, multi-modal AU+gaze fusion~\cite{abedi2021affect}) push
the supervised ceiling to $60$--$73\%$ accuracy, but no published
work investigates the frozen-feature/zero-shot regime in depth. Our
work is complementary: rather than seek a higher supervised number,
we characterise the regime that practitioners and zero-shot
deployers actually encounter.

\paragraph{Vision-language models as visual representation
backbones.} CLIP~\cite{radford2021clip} and its
successors---SigLIP~\cite{zhai2023siglip},
DINOv2~\cite{oquab2023dinov2}, TIPSv2~\cite{cao2026tipsv2}---are
routinely used as frozen backbones for downstream visual recognition
\cite{wortsman2022robust,kumar2022finetuning}. Affect recognition
benchmarks (FER2013~\cite{goodfellow2013fer2013},
AffectNet~\cite{mollahosseini2017affectnet}) have been shown to
benefit substantially from frozen CLIP features. We confirm this
finding on FER2013 ($\kappa \approx 0.55$) and use it as a
cross-task control bounding the scope of our DAiSEE claims.

\paragraph{Linear probing and the role of identity.} Belrose
~\etal~\cite{belrose2023leace} introduce Linear Erasure of Affine
Concepts (LEACE), an optimal linear projection that removes any
linear classifier of a concept while minimising distortion.
\cite{ravfogel2020rlace} and~\cite{ravfogel2022linear} provide
related iterative null-space projection methods (INLP). Both lines
of work show that linear erasure is sufficient for some semantic
concepts but can destroy correlated downstream signals. We apply
LEACE to DAiSEE and find the engagement-identity entanglement is
strong enough that linear erasure of identity also erases
engagement, suggesting these methods are not a remedy for the
DAiSEE-style affective benchmarks.

\paragraph{Subject-invariant affect models and personalization.}
Subject-invariant emotion recognition has a long tradition: DANN
gradient-reversal heads~\cite{ganin2015dann}, person-invariant
adversarial features~\cite{li2017cross}, supervised contrastive
losses with subject-stratified negatives~\cite{khosla2020supcon}.
Personalization at inference time has been studied in BCI
\cite{lotte2018review} and affective computing~\cite{rudovic2018personalized}
with success when a labelled support set per subject is available.
We instead consider \emph{label-free} personalization---a single
self-supervised neutral-anchor frame per subject---and show that
in the linear regime it inherits the same identity-engagement
entanglement that makes adversarial methods unstable.

\paragraph{Ordinal classification.} For ordinal labels, $\kappa_q$
is sensitive to the choice of decision boundaries. Frank \&\ Hall
\cite{frank2001ordinal} and CORN~\cite{shi2023corn} provide
threshold-based ordinal classifiers; Menon~\etal~\cite{menon2021logit}
provide post-hoc logit-adjustment for class imbalance. We combine
class-balanced LR with val-tuned ordinal thresholds (the
expected-class~$\mathbb{E}[y]$ score binned at three thresholds);
this single intervention is responsible for a $+0.025$ absolute
$\kappa_q$ lift over argmax decoding---larger than any other
intervention we tested.
